% Part: normal-modal-logic
% Chapter: completeness
% Section: introduction

\documentclass[../../../include/open-logic-section]{subfiles}

\begin{document}

\olfileid[hi]{nml}{com}{int}

\olsection{परिचय}

यदि $\Sigma$ कोई मोडल तंत्र है, तो यथार्थता प्रमेय स्थापित करता है कि यदि
$\Sigma \Proves !A$, तो $!A$ प्रतिरूपों के ऐसे किसी भी वर्ग
$\mClass{C}$ में मान्य है जिसमें $\Sigma$ के सभी !!{formula}ों के सभी
दृष्टांत मान्य हैं। विशेषतः इसका अर्थ है कि यदि $\Log{K} \Proves !A$,
तो $!A$ सभी प्रतिरूपों में सत्य है; यदि $\Log{KT} \Proves !A$, तो $!A$
सभी स्वपरावर्ती प्रतिरूपों में सत्य है; यदि $\Log{KD} \Proves !A$, तो
$!A$ सभी सीरियल प्रतिरूपों में सत्य है; इत्यादि।

पूर्णता, यथार्थता का विलोम है: उदाहरणार्थ, \Log{K} के पूर्ण होने का अर्थ
है कि यदि कोई !!a{formula}~$!A$ मान्य है, तो $\Proves !A$। पूर्णता सिद्ध
करना यथार्थता सिद्ध करने से कहीं कठिन है। पहले प्रतिधनात्मक रूप पर विचार
करना उपयोगी है: \Log{K} तभी पूर्ण है जब $\Proves/ !A$ होने पर कोई
प्रतिप्रतिरूप, अर्थात् ऐसा प्रतिरूप~$\mModel{M}$ हो कि
$\mSat/{M}{!A}$। तुल्यतः ($!A$ का निषेध लेकर), हम सिद्ध कर सकते हैं कि जब
भी $\Proves/ \lnot !A$, तब $!A$ का कोई प्रतिरूप है। ऐसा प्रतिरूप बनाते
समय हम $!A$ में निहित सूचना का उपयोग कर सकते हैं। विशिष्ट !!{formula}ों के
लिए प्रतिरूप खोजते समय हम प्रायः यही करते हैं: जैसे, यदि हमें $p \lif
\Box q$ का प्रतिप्रतिरूप खोजना है, तो हम जानते हैं कि उसमें ऐसा जगत होना
चाहिए जहाँ $p$ सत्य और $\Box q$ असत्य हो। और जिस जगत पर $\Box q$ असत्य
है, उससे अभिगम्य ऐसा जगत होना चाहिए जहाँ $q$ असत्य हो। हमें बस इतना ही
जानना है: कौन-से जगत !!{propositional variable}ों को सत्य बनाते हैं और कौन-से
जगत किन जगतों से अभिगम्य हैं।

किंतु पूर्णता सिद्ध करते समय हमारे पास ऐसा कोई विशिष्ट !!{formula}~$!A$
नहीं है जिसके लिए हम प्रतिरूप बना रहे हों। हम स्थापित करना चाहते हैं कि
हर ऐसे $!A$ के लिए प्रतिरूप विद्यमान है जिसके लिए $\Proves/[\Sigma]
\lnot !A$। यह न्यूनतम अपेक्षा है, क्योंकि \emph{यदि}
$\Proves[\Sigma] \lnot !A$, तो यथार्थता के अनुसार $!A$ का कोई ऐसा
प्रतिरूप नहीं है जिसमें $\Sigma$ सत्य हो। अब ध्यान दें कि
$\Proves/[\Sigma] \lnot !A$ तभी जब $!A$, $\Sigma$-संगत है। (स्मरण रहे
कि $\Sigma \Proves/[\Sigma] \lnot !A$ और $!A \Proves/[\Sigma]
\lfalse$ तुल्य हैं।) अतः हमारा कार्य प्रत्येक $\Sigma$-संगत !!{formula}
के लिए प्रतिरूप बनाना है।

हम जिस युक्ति का उपयोग करेंगे, वह !!{formula}ों का ऐसा $\Sigma$-संगत
समुच्चय खोजना है जिसमें $!A$ के साथ वे अन्य सूत्र भी हों जो हमें बताएँ
कि $!A$ को सत्य बनाने वाला जगत कैसा होना चाहिए। ऐसे समुच्चय
\emph{पूर्ण} $\Sigma$-संगत समुच्चय हैं। $!A$ को सत्य बनाने के लिए केवल
एक जगत वाला प्रतिरूप बनाना पर्याप्त नहीं; उसमें अनेक जगत और एक अभिगम्यता
संबंध होना आवश्यक होगा। $!A$ वाला पूर्ण $\Sigma$-संगत समुच्चय
$\Box !B$ और~$\Diamond !C$ रूप के अन्य !!{formula} भी रखेगा। सभी अभिगम्य
जगतों में $!B$ सत्य होना चाहिए; कम-से-कम एक में $!C$ सत्य होना चाहिए। इसे
पूरा करने के लिए हम जगतों के समुच्चय के आधार के रूप में सभी संभावित पूर्ण
$\Sigma$-संगत समुच्चय ले लेंगे। कठिन बात यह तय करना होगी कि हमारे
प्रतिरूप में एक पूर्ण $\Sigma$-संगत समुच्चय कब दूसरे से अभिगम्य माना जाए।

हम दिखाएँगे कि इस प्रकार परिभाषित प्रतिरूप में $!A$ किसी जगत पर---जो स्वयं
भी पूर्ण $\Sigma$-संगत समुच्चय है---तभी सत्य है जब $!A$ उस समुच्चय का
!!a{element} हो। यदि $!A$, $\Sigma$-संगत है, तो वह कम-से-कम एक पूर्ण
$\Sigma$-संगत समुच्चय का !!a{element} होगा (यह तथ्य हम सिद्ध करेंगे), और
इसलिए कोई ऐसा जगत होगा जहाँ $!A$ सत्य है। अतः हमारे पास ऐसा एक प्रतिरूप
होगा जिसमें प्रत्येक $\Sigma$-संगत !!{formula}~$!A$ किसी जगत पर सत्य है।
यही एक प्रतिरूप $\Sigma$ का \emph{विहित} प्रतिरूप है।

\end{document}
